La frecuencia es el número de ciclos completos de la onda por segundo, medida en herzios (\textit{Hz}) y determina el tono de lo que escuchamos. Las notas tienen sus propias frecuencias, por ejemplo \textit{440Hz} corresponde a \textit{La4}, la voz humana suele ir en el rango 128-256 Hz , los ultrasonidos se consideran los superiores a 20 kHz y los infrasonidos los inferiores a 20 Hz \cite{libretexts_sound}. El rango de audición humana oscila entre los 20Hz y los 20kHz.

Por otro lado , el tono está relacionado con la \textit{frecuencia fundamental} ( aunque en sonidos complejos intervienen también armónicos, que son múltiplos de la frecuencia fundamental , lo cual afecta también al timbre , siendo el timbre lo que nos permite distinguir una misma nota tocada en una guitarra, un violín, o una flauta. Este tipo de frecuencias suelen extenderse por todo el rango frecuencial audible, por lo que , en relación con el proyecto, su presencia suele dificultar el trabajo de identificación de elementos sonoros.)